\documentclass[10pt,a4paper]{amsart}
\usepackage[margin=19mm]{geometry}
\usepackage{amsmath,amssymb,mathtools}
\usepackage[scr=rsfs,cal=euler]{mathalfa}
\usepackage{xfrac}
\usepackage[unicode,hidelinks,bookmarks=false]{hyperref}
\newcommand{\ie}{i.e.\ }
\newcommand{\cf}{cf.\ }
\newcommand{\resp}{respectively\ }
\newcommand{\Subsection}{\subsection}
\providecommand{\nobreakdash}{\nobreak\hbox{-}}
\setlength{\emergencystretch}{2em}
\multlinegap0pt
\usepackage{mathtools}
\usepackage{amssymb}
\usepackage{xcolor}
\usepackage{tikz}
\usepackage[shortlabels]{enumitem}
\newtheorem{proposition}{Proposition}
\newcommand{\R}{\mathbb{R}}
\def\cd{\centerdot}
\def\d{\partial}
\newcommand{\dd}{\mathrm{d}}
\def\ind{\mathbf 1}
\def\ovl{\overline}
\DeclareMathOperator{\supp}{supp}
\def\ve{\varepsilon}
\title{The initial-to-final-state inverse problem with critically-singular potentials}
\author{Manuel Ca\~nizares, Pedro Caro, Ioannis Parissis, Thanasis Zacharopoulos}
\date{Source: arXiv:2602.12122v1 (CC BY 4.0)}
\begin{document}
\maketitle

\noindent\emph{Source and licence.} The initial-to-final-state inverse problem with critically-singular potentials. Version arXiv:2602.12122v1 (CC BY 4.0), distributed by arXiv under the Creative Commons Attribution 4.0 International licence, http://creativecommons.org/licenses/by/4.0/, as recorded in the arXiv licence metadata for this exact version and shown on the abstract page https://arxiv.org/abs/2602.12122v1. Full licence notice: \texttt{licenses/arxiv-2602.12122-CC-BY-4.0.txt}.\par\smallskip

\noindent\textit{Editorial scope.} Counted unit: the article's proposition prp:integration\_by\_parts\_unbounded\_no-endpoint (source0.tex lines 670--682). The article's complete original proof of it is delivered with it (source0.tex lines 710--865, one \texttt{proof} environment of 156 lines, which the article places away from the statement). No mathematics is written, repaired or re-derived: the statement and the proof are the article's own lines, in the article's own notation, and no retained line is substituted or re-flowed.  Retained uncounted as the source-local context the proof needs: lines 688--705.  Nothing was omitted from any retained span, so the package carries no omission record.  No retained line contains a citation, so the wrapper carries no bibliography and no bibliography entry.  No retained line contains a figure environment, an image-inclusion command or drawing code, so no image is delivered and the package declares no asset dependency.  Editorial formatting only, in addition: an AMS \texttt{amsart} preamble that restates the article's own declarations and macros used by the retained lines, namely the environments \texttt{proposition} on separate counters, together with the standard packages those lines need.  The retained environments print in this document as Proposition 1 in order of appearance, whereas the article prints the same items with the numbering of its own full document.  The complete native source of the article as distributed in the arXiv e-print for this exact version is bundled unchanged as \texttt{sources/194522/source0.tex}.
\begin{proposition}\label{cor:integration_by_parts_unbounded_no-endopint2}\sl
Let $n\geq 2$ and $(r,p)$ be a Strichartz pair. Let
\[
v \in C\left([0,T];L^p(\R^n)\right)
\]
and
\[
u \in C\left([0,T];L^2(\R^n)\right) \cap L^r\left((0,T);L^p(\R^n)\right)\quad\text{with}\quad u(0,\cd)=u(T,\cd)=0,
\]
be such that
\[
\left(i\partial_t+\Delta\right)u, \left(i\partial_t+\Delta\right)v \in L^{r'}\big((0,T);L^{p'}(\R^n)\big).
\]
Then
\[
\int_\Sigma \left(i\partial_t+\Delta\right)u \ovl v = \int_\Sigma u \ovl{\left(i\partial_t+\Delta\right)v}.
\]
\end{proposition}

\begin{proposition}\label{prp:integration_by_parts_unbounded_no-endpoint}\sl Let $n\geq 2$ and $(r,p)$ be a Strichartz pair. Assume
\[
u,v \in C\left ([0,T];L^2(\R^n)\right) \cap L^r\left((0,T);L^p(\R^n)\right)
\]
be such that
\[
(i\partial_t+\Delta)u, (i\partial_t+\Delta)v \in L^{r'}\big((0,T);L^{p'}(\R^n)\big).
\]
Then
\[
\int_{\Sigma} \left[(i\partial_t+\Delta)u \ovl v - u \ovl{(i\partial_t+\Delta)v}\right] = i \int_{\R^n} \left[u(T,\cd)\ovl{v(T,\cd)}-u(0,\cd)\ovl{v(0,\cd)}\right].
\]
\end{proposition}

\begin{proof} If we additionally assume that $v\in C([0,T];L^2(\R^n))$, the conclusion follows directly from Proposition~\ref{prp:integration_by_parts_unbounded_no-endpoint}, since the boundary terms vanish due to the conditions $u(0,\cd)=u(T,\cd)=0$. The main part of the proof is therefore devoted to removing the assumption $v\in C([0,T];L^2(\R^n))$ and reducing the case $v\in C([0,T];L^p(\R^n))$ to the $L^2$-case via approximation.

Let $\chi$ be a smooth bump function on $\R^n$ such that $0\leq \chi \leq 1$, $\chi\equiv 1$ on $\{|x|\leq 1\}$, and $\supp \chi \subset \{|x|\leq 2\}$. Let also $\phi\in \mathcal S(\R^n)$ be a nonnegative bump function with $\supp \phi \subset B(0,1)$, satisfying $\int_{\R^n}\phi=1$. For $\varepsilon>0$, let $R=R(\varepsilon)>0$ be a function such that
\[
\lim_{\varepsilon\to0} R(\varepsilon)=+\infty,
\]
which will be chosen in the course of the proof.

We define the rescaled bump functions $\chi^R(x)\coloneqq \chi(x/R)$ and $\phi_\varepsilon(x)\coloneqq \varepsilon^{-n}\phi(x/\varepsilon)$ and for general $v \in C([0,T];L^p(\R^n))$, we set
\[
v_{R,\varepsilon} \coloneqq \chi^R (\phi_\varepsilon * v).
\]
Here the convolution $\phi_\varepsilon * v$ is taken only in the space variables.  Since $v\in C([0,T];L^p(\R^n))$ with $p\geq 2$ and $\chi^R$ has compact support, it follows that
\[
v_{R,\varepsilon}\in C\left([0,T];L^2(\R^n)\right) \cap L^r\left((0,T);L^p(\R^n)\right).
\]
Now an application of the Leibniz rule in the sense of distributions yields
\begin{equation}\label{eq:Leibniz_distributional}
\left(i \d_t + \Delta \right)v_{R, \ve} = \chi^R \left(i\d_t + \Delta\right)(v*\phi_\ve) +  ( v*\phi_\ve ) \Delta(\chi^R)  + 2 ( \nabla \chi^R ) \cdot \nabla (v*\phi_\ve),
\end{equation}
in $\mathcal{D}' (\Sigma)$. We record for later use the easy identities
\begin{equation}\label{eq:scalederiv}
\begin{split}
& \Delta(\chi^R) = R^{-2} (\Delta \chi)^R \ind_{\{R\leq |x| \leq 2R\}}, \qquad \nabla (\chi^R) = R^{-1} (\nabla \chi)^R \ind_{\{R\leq |x| \leq 2R\}}, 
\\
&  \nabla (\phi_\ve) = \ve^{-1} (\nabla \phi)_{\ve}. 
\end{split}
\end{equation}

We estimate each term in \eqref{eq:Leibniz_distributional}. First, we have
\[
\chi^R (i\d_t + \Delta)(v*\phi_\ve) = \chi^R \phi_\ve * (i\d_t + \Delta)v \in L^{r'}\big((0, T); L^{p'}(\R^n)\big).
\]
Moreover, the function
\[
( v*\phi_\ve ) \Delta(\chi^R) = R^{-2} (\Delta \chi)^R (v*\phi_\ve)
\]
has support contained in $\{x\in\R^n:\, R\leq |x|\leq 2R\}$ by the definition of $\chi$. Since $p\geq 2$ we have that $p'\leq 2 \leq p$, and hence $(v*\phi_\ve) \Delta (\chi^R)\in C([0, T]; L^{p'}(\R^n))$.  Similarly, we also infer that 
\[
( \nabla \chi^R ) \nabla (v*\phi_\ve) = R^{-1} \ve^{-1} (\nabla \chi)^R (v* (\nabla \phi)_\ve) \in C\big ([0, T]; L^{p'}(\R^n)\big). 
\]
Therefore, $(i \d_t + \Delta )v_{R, \ve} \in L^{r'}((0, T); L^{p'}(\R^n))$ and we can use the integration by parts formula of Proposition~\ref{prp:integration_by_parts_unbounded_no-endpoint} with the right hand side equal to zero to get 
\[
\int_\Sigma [(i\d_t + \Delta)u] \ovl{v_{R,\ve}} = \int_\Sigma \ovl{(i\d_t + \Delta) v_{R, \ve}}  u.
\]
Notice that, in order to complete the proof it will be enough to show that 
\begin{equation}\label{eq: Cor_int_by_parts_approx1}
\lim_{\ve \to 0} \int_\Sigma \left[(i\d_t + \Delta) u\right] \ovl{(v_{R, \ve} - v)} = 0, 
\end{equation}
and
\begin{equation}\label{eq: Cor_int_by_parts_approx2}
\lim_{\ve \to 0} \int_\Sigma u  \ovl{\left(i\d_t + \Delta\right)(v_{R, \ve} - v)} = 0.
\end{equation}

We begin with the proof of~\eqref{eq: Cor_int_by_parts_approx1}. We have that 
\begin{align*}
\Big| \int_\Sigma  [(i\d_t + \Delta) u] \ovl{(v_{R, \ve} - v)} \Big| &\leq 
\| (i\d_t + \Delta)u \|_{L^{r'}((0, T); L^{p'}(\R^n))} \| v_{R, \ve} - v \|_{L^r((0, T); L^p(\R^n))} \\
&\leq \| (i\d_t + \Delta)u \|_{L^{r'}((0, T); L^{p'}(\R^n))} \| (\chi^R - 1) v\|_{L^r((0, T); L^p(\R^n))} \\
&+ \| (i\d_t + \Delta)u \|_{L^{r'}((0, T); L^{p'}(\R^n))} \| \phi_\ve * v - v \|_{L^r((0, T); L^p(\R^n))}
\end{align*}
which implies~\eqref{eq: Cor_int_by_parts_approx1} as $\ve \to 0$ by dominated convergence. 

We now turn to the proof of the more challenging limit~\eqref{eq: Cor_int_by_parts_approx2}. We have that 
\begin{equation}\label{eq: Cor_int_by_parts_approx2_split} 
\begin{split}
\left| \int_\Sigma \big[ \ovl{ (i\d_t + \Delta)(v_{R,\ve} -v) } \big] u \right| &\leq
\left| \int_\Sigma \big[ \ovl{ \chi^R [(i\d_t+\Delta)v]*\phi_\ve - (i\d_t+\Delta)v }\big] u \right| 
\\
&\qquad +\left| \int_\Sigma \overline{\Delta(\chi^R) (v*\phi_\ve)} u \right| 
+ 2 \left| \int_\Sigma \overline{\nabla (\chi^R) \cdot \nabla (v*\phi_\ve)} u \right|
\\
&\eqqcolon \mathrm{I}(\ve,R)+\mathrm{II}(\ve,R)+\mathrm{III}(\ve,R).
\end{split}  
\end{equation}

We begin by estimating $\mathrm I(\ve,R)$. We have
\[
\begin{split}
\mathrm{I}(\ve,R) &\leq  \int_\Sigma \left| \left[\ovl{\phi_\ve*\left(i\d_t+\Delta\right)v - \left(i\d_t+\Delta\right)v } \right] \chi^R \right| |u|  + \int_\Sigma \left|\chi^R - 1\right| \left|\left(i\d_t +\Delta\right)v \right| |u| 
\\
&\leq \left\| \phi_\ve*\left(i\d_t+\Delta\right)v - \left(i\d_t+\Delta\right)v \right\|_{L^{r'}((0, T); L^{p'}(\R^n))} \| u \|_{L^r((0, T); L^p(\R^n))} 
\\
& \qquad + \int_\Sigma \left|\chi^R - 1\right| \left|\left(i\d_t +\Delta\right)v \right| |u|.
\end{split}
\]
In the right-hand side of the display above, the first summand tends to $0$ as $\varepsilon\to0$ by the convergence of the approximate identity $\phi_\varepsilon*(i\d_t+\Delta)v$ in $L^{p'}(\R^n)$, while the second summand tends to $0$ as $R\to\infty$ by dominated convergence. Consequently, $\mathrm I(\ve,R(\ve))\to 0$ as $\varepsilon\to0$ with any choice of $R=R(\varepsilon)\to\infty$ as $\ve \to 0$.

We estimate the worst term $\mathrm{III}(\varepsilon,R)$. By the assumptions on $\chi,\phi$ and \eqref{eq:scalederiv}, we have
\[
\left|\nabla(\chi^R)\cdot\nabla(v*\phi_\varepsilon)\right|= \ve^{-1} \left|\nabla(\chi^R)\cdot (v*(\nabla\phi)_\varepsilon)\right| \lesssim R^{-1}\varepsilon^{-1} \ind_{\{R\le |x|\le 2R\}}  \left(|v|*|(\nabla\phi)_\varepsilon|\right).
\]
Since $\supp (\nabla\phi)_\varepsilon \subset B(0,\varepsilon)$, it follows that, provided
$\varepsilon\ll R$,
\[
\ind_{\{R\le |x|\le 2R\}}\left(|v|*|(\nabla\phi)_\varepsilon|\right) \leq \left(\ind_{\{|x|\simeq R\}}|v|\right)*|(\nabla\phi)_\varepsilon|.
\]
Therefore,
\[
\begin{split}
\mathrm{III}(\varepsilon,R) &\lesssim R^{-1}\varepsilon^{-1}\int_\Sigma \left[\left(\ind_{\{|x|\simeq R\}}|v|\right)*|(\nabla\phi)_\varepsilon|\right](t,x) |u(t,x)|\,\dd x\,\dd t
\\
&\le R^{-1}\varepsilon^{-1}  \left\|\left(\ind_{\{|x|\simeq R\}}|v|\right)*|(\nabla\phi)_\varepsilon|\right\|_{L^1((0,T);L^2(\R^n))} \|u\|_{C([0,T];L^2(\R^n))}
\\
&\le R^{-1}\varepsilon^{-1} \|v\|_{L^1((0,T);L^2(\{|x|\simeq R\}))}\|u\|_{C([0,T];L^2(\R^n))} 
\\
&\lesssim R^{-1}\varepsilon^{-1}  R^{n(\frac12-\frac1p)} \|v\|_{L^1((0,T);L^p(\{|x|\simeq R\}))}\|u\|_{C([0,T];L^2(\R^n))} .
\end{split}
\]
Here, we used Cauchy--Schwarz in $x$ and H\"older in $t$ to pass to the second line, Young's inequality in $x$ in passing to the third line, together with $\|(\nabla\phi)_\varepsilon\|_{L^1(\R^n)}\lesssim 1$, and finally H\"older on the annulus $\{|x|\simeq R\}$ to pass from $L^2(\{|x|\simeq R\})$ to $L^p(\{|x|\simeq R\})$ with $p\geq2$. Using the calculation
\[
-1+n\left(\frac12-\frac1p\right)= n\left(\frac1{p_n}-\frac1p\right), \qquad  \frac1{p_n}=\frac12-\frac1n,
\]
and writing $R=R(\varepsilon)$, we have proved
\[
\mathrm{III}(\varepsilon,R(\varepsilon)) \lesssim \varepsilon^{-1}R(\varepsilon)^{-n(\frac1p-\frac1{p_n})} \|v\|_{L^1((0,T);L^p(\{|x|\simeq R(\varepsilon)\}))} \|u\|_{C([0,T];L^2(\R^n))}.
\]

If $p<p_n$, then $n(\frac1p-\frac1{p_n})>0$, and choosing
\[
R(\varepsilon)\coloneqq \varepsilon^{-\frac{2}{n(\frac1p-\frac1{p_n})}}
\]
gives $\mathrm{III}(\varepsilon,R(\varepsilon))\lesssim \varepsilon$, hence $\mathrm{III}(\varepsilon,R(\varepsilon))\to0$ as $\varepsilon\to0$.

In the endpoint case $p=p_n$ we have to be more careful since we get that
\[
\mathrm{III}(\ve,R(\ve))
\leq \frac{1}{\ve}  \|u\|_{C([0, T]; L^2(\R^n))}  \|v\|_{L^1\left((0, T); L^{ p_n}(|x|\simeq R(\ve))\right)} .
\]
Notice that
\[
v\ind_{\{|x|\simeq R\}}\ind _{(0,T)} \in C\left([0,T];L^{p_n}(\R^n)\right) \subset L^1\left((0,T);L^{p_n}(\R^n)\right),
\]
and that
\[
v(t,x)\ind_{\{|x|\simeq R\}}\ind_{(0,T)}(t)\to0 \quad\text{as }R\to\infty
\]
for every $(t,x)\in\Sigma$. By dominated convergence,
\[
\lim_{R\to\infty} \|v\|_{L^1((0,T);L^{p_n}(\{|x|\simeq R\}))}=0.
\]
Hence, for every $\ve>0$ there exists $R(\ve)>0$ which can be chosen to also satisfy $R(\ve)>\ve^{-1}$ such that $\|  v \|_{L^1((0,T); L^{p_n}(\{|x|\simeq R\}))}<\ve ^2$ for all $R\geq R(\ve)$. This shows that in the endpoint case $p=p_n$ one can choose $R(\ve)$, with $\lim_{\ve\to0 }R(\ve)=+\infty$, so that
\[
\mathrm{III}(\ve,R(\ve)) \lesssim \frac{1}{\ve}  \|u\|_{C([0, T]; L^2(\R^n))}  \|v\|_{L^1((0, T); L^{p_n}(|x|\simeq R(\ve)))} \to 0 \quad \text{as}\quad \ve\to 0.
\]

The term $\mathrm{II}(\varepsilon,R(\varepsilon))$ is estimated similarly. In particular, in the endpoint case $p=p_n$ we obtain
\[
\mathrm{II}(\ve,R(\ve)) \lesssim  \frac{1}{R(\ve)} \|u\|_{C([0, T]; L^2(\R^n))} \|v\|_{L^1((0, T); L^{p_n}(\R^n))} \to 0 \quad \textup{as} \quad \ve\to 0. 
\]
Combining the estimates we see that there is a choice $R(\ve)\to \infty$ when $\ve\to 0$ such that 
\[
\lim_{\ve\to 0} \left[\mathrm{I}(\ve,R(\ve))+\mathrm{II}(\ve,R(\ve))+\mathrm{III}(\ve,R(\ve))\right]=0
\]
so the right hand side of~\eqref{eq: Cor_int_by_parts_approx2_split} tends to 0 as $\ve\to 0$. This  shows~\eqref{eq: Cor_int_by_parts_approx2} and completes the proof of  Proposition~\ref{cor:integration_by_parts_unbounded_no-endopint2}. 
\end{proof}

\end{document}
